\section{서론}
\label{sec:intro}

종단간(end-to-end, E2E) 자율주행은 센서 관측과 자차 상태로부터 미래 궤적이나 제어 의도를 직접 예측한다 \cite{hu2023uniad,jiang2023vad}. 최근의 추론 보강형 E2E/VLA 연구는 여기에 행동 이유를 설명하는 Chain-of-Causation(CoC)을 결합한다 \cite{sima2024drivelm,tian2024drivevlm,wang2025alpamayo}. CoC는 ``전방 보행자 때문에 정지한다''처럼 상황 인식, 요구 행동과 근거를 함께 제시하므로 학습 자료의 의미를 풍부하게 하고 모델 출력을 검토하는 단서가 될 수 있다.

그러나 CoC의 품질을 각 사건에서만 검사하면 중요한 오류를 놓친다. 첫 사건이 ``보행자가 통과할 때까지 정지한다''라고 요구하고 다음 사건이 ``가속하여 진행한다''라고 요구한다고 하자. 두 문장은 각각 현재 행동을 설명하는 형식으로는 자연스러울 수 있다. 하지만 보행자의 통과가 확인되지 않았다면 첫 문장의 정지 의무는 여전히 활성이고, 두 번째 행동은 그 의무와 충돌한다. 반대로 해제 조건이 충족됐는데도 정지 의무가 남아 있으면 이후 진행을 계속 막는다. 이러한 오류는 현재 문장만이 아니라 이전 사건에서 시작된 의무와 그 해제 이력을 보아야 드러난다.

\begin{figure*}[t]
\centering
\resizebox{0.95\textwidth}{!}{%
\begin{tikzpicture}[
  main/.style={draw, rounded corners=2pt, align=center, minimum height=12mm,
    text width=3.05cm, font=\footnotesize},
  stage/.style={draw, rounded corners=2pt, align=center, minimum height=8mm,
    text width=3.45cm, font=\scriptsize, fill=black!3},
  data/.style={main, fill=observed!10},
  reason/.style={main, fill=derived!10},
  contract/.style={main, fill=assumed!12},
  route/.style={main, fill=neutral!8},
  flow/.style={-{Latex[length=2mm]}, thick, draw=black!80},
  guide/.style={-{Latex[length=1.5mm]}, thin, dashed, draw=black!45}
]

% Main narrative.
\node[data] (log) at (0,0) {Recorded driving data\\ego trajectory};
\node[reason] (llm) at (4.0,0) {Sequential VLM/LLM\\CoC annotations};
\node[contract] (obl) at (8.0,0) {Persistent obligation\\release and action order};
\node[contract] (monitor) at (12.0,0) {Stateful Python checker\\UPPAAL observer};
\node[route] (route) at (16.0,0) {Separated reports\\natural sample\\synthetic regression\\response/provenance analyses};

\draw[flow] (log) -- coordinate[midway] (m1) (llm);
\draw[flow] (llm) -- coordinate[midway] (m2) (obl);
\draw[flow] (obl) -- coordinate[midway] (m3) (monitor);
\draw[flow] (monitor) -- coordinate[midway] (m4) (route);

% Process interpretation row.
\node[stage] (s1) at (2.0,-1.55) {Post-hoc labeling from\\recorded clips};
\node[stage] (s2) at (6.0,-1.55) {Within-scene ordering\\contract compilation};
\node[stage] (s3) at (10.0,-1.55) {Three-valued evidence\\obligation-state update};
\node[stage] (s4) at (14.0,-1.55) {Verdict and first\\violating event};

\draw[guide] (m1) -- (s1.north);
\draw[guide] (m2) -- (s2.north);
\draw[guide] (m3) -- (s3.north);
\draw[guide] (m4) -- (s4.north);

% Takeaway.
\node[draw=assumed, rounded corners=2pt, align=center, font=\footnotesize,
  fill=assumed!8, text width=14.0cm] (key) at (8.0,-3.0)
  {Key point: stateful rules can retain obligations across events, but synthetic rule coverage is not evidence of detection accuracy on unseen natural CoCs.};

\end{tikzpicture}%
}

\bifigcaption{추론 보강형 E2E 주행 데이터의 연속 CoC를 상태 계약으로 변환하는 전체 흐름. 지속 의무, 해제 조건, 행동 순서와 삼값 증거를 명시하고, 자연 표본 검토와 합성 회귀 시험 및 보조 반응ㆍ출처 분석의 결과를 분리한다.}{Overall pipeline for compiling sequential CoCs into stateful contracts while separating natural-sample review, synthetic regression, and auxiliary response/provenance analyses.}
\label{fig:overview-pipeline}
\end{figure*}

그림~\ref{fig:overview-pipeline}은 기록된 영상ㆍ자차 궤적과 사후 생성 CoC로부터 검증 결과를 얻는 전체 흐름을 요약한다. 본 연구는 CoC를 바로 제어 명령이나 안전 판정으로 사용하지 않고, 명시된 계약 의미론에 따라 지속 의무와 해제 조건을 가진 사건열로 변환한다. Python 검사기는 사건열의 상태를 갱신하고, 같은 전이 규칙에서 생성한 UPPAAL 관찰 모델은 정식 시험 사례에서 위반 상태와 최초 위반 사건을 재현한다. 자연 표본의 사람 검토, 합성 회귀 시험과 보조 반응ㆍ출처 결과는 하나의 성능 수치로 합치지 않는다.

본 논문이 푸는 문제는 \emph{연속 CoC가 선언한 규칙 아래에서 순서ㆍ상태 면에서 자기 일관적인가}이다. CoC를 단순 행동 레이블이 아니라 지속 의무, 해제 조건과 필수 행동 순서를 갖는 행동 계약으로 해석한다. 텍스트나 관측 자료가 해제 상태를 말하지 않으면 참이나 거짓으로 추정하지 않고 {\unknown}으로 유지한다. 상태형 Python 검사기는 같은 장면의 사건을 시간순으로 읽으며 활성 의무를 추적한다. 별도의 \uppaal 관찰 모델은 같은 전이 규칙을 구현하여 정해진 시험 사례에서 판정 일치와 최초 충돌 상태를 확인한다. 두 구현은 공통 의미론을 공유하므로 이 일치는 독립적인 의미 타당성 검증이 아니다.

이 검증은 두 맥락에서 활용 가능성을 갖는다. 첫째, 연속 추론 자료를 학습에 사용하기 전에 선언한 계약과 충돌하는 행동 감독 후보를 표시할 수 있다. 둘째, 활성 의무와 위반 경로를 통해 어떤 과거 조건이 후속 행동과 충돌했는지 추적할 수 있다. 다만 이런 활용 효과는 자연 자료의 실제 오류나 학습 성능으로 검증하지 않았으며, 본 실험은 \emph{과거 의무 상태를 유지하는 품질 점검 규칙의 실행 가능성}을 평가한다.

본 논문의 중심 연구 질문은 다음 하나이다.

\begin{quote}
\textbf{연구 질문:} 미리 선언한 CoC 계약 의미론에서 상태 기반 검사는 사건 간 상태가 필요한 합성 모순 유형을 사건별 검사와 구분하여 판정할 수 있으며, 같은 의미론의 UPPAAL 구현은 그 판정과 최초 위반 상태를 재현할 수 있는가?
\end{quote}

자연 자료와 합성 시험의 역할은 분리한다. 자차 궤적이 있는 다중-CoC 90장면에서 309개 장면 내부 인접 창을 추출하고, 그중 선택한 27개 독립 장면을 네 검토자의 블라인드 판정과 합의로 탐색한다. 실행 규칙은 자연 문장을 변형한 평가셋이 아니라 네 모순 유형별 10쌍, 총 40개의 결정적 합성 중간표현 기준--변형 쌍으로 점검한다. 생성기는 상태형 검사 결과를 사후조건으로 강제하므로, 사건별 기준선과의 차이는 상태 기억이 필요한 규칙의 구조적 차이와 시험 커버리지를 보여 줄 뿐 미관측 자연 오류에 대한 검출 성능을 추정하지 않는다. 실제 \texttt{verifyta} 실행은 일곱 정식 시험 사례에서 Python과 UPPAAL 구현의 판정 일치를 확인한다.

본 논문의 핵심 기여는 세 가지이다.

\begin{enumerate}
    \item 연속 CoC를 지속 의무, 명시적 해제 조건, 행동 순서와 삼값 증거로 표현하는 공통 계약 중간표현을 제시한다.
    \item 상태형 Python 검사기와 같은 전이 의미론의 UPPAAL 관찰 모델을 구현하여, 판정과 최초 위반 상태를 추적할 수 있게 한다.
    \item 선택한 자연 장면의 사람 검토 결과와 결정적 합성 회귀 시험을 분리해 보고한다. 합성 기준--변형 40쌍에서는 상태형 규칙이 40/40, 사건별 규칙이 20/40 변형을 모순으로 판정했고, 일곱 UPPAAL 시험 사례의 판정 불일치는 0/7이었다.
\end{enumerate}

기존의 단일 사건 시간 반응, 반복 정책 민감도와 장면 출처 분석은 계약 입력과 경계 정책을 설명하는 보조 근거로 유지하되 중심 합성 회귀 시험과 분리한다. 본 논문의 주장은 선언한 CoC 텍스트 계약 의미론의 실행 가능성과 추적성에 한정된다. 자연 자료에서의 오류 검출률, 실제 교통법규 위반, 충돌 위험, 물리적 차량 안전 또는 학습 성능 향상을 주장하지 않는다.
